\section*{Conclusion du chapitre}
\addcontentsline{toc}{section}{Conclusion du chapitre}

Ce chapitre a été dédié à l'étude et au développement de méthodologies de diagnostic et de pronostic appliquées aux composants de stockage et de conversion d'énergie d'un micro-réseau.

Dans un premier temps, des approches de diagnostic ont été développées pour les batteries lithium-ion, pour une estimation robuste du SoH. Une campagne expérimentale suivant l'évolution des caractéristiques statiques et dynamiques des batteries a été menée, et des modèles de vieillissement en ont été tirés pour une évaluation de l'état de santé.

Parallèlement, une part égale de l'étude a été consacrée aux systèmes à hydrogène, ie. la pile à combustible et l'électrolyseur. Des
modèles de vieillissement des composants ont été introduits et c'est ensuite le filtre de Kalman étendu qui a été retenu puis appliqué pour suivre l'évolution de l'état de santé des composants en fonction des points de fonctionnement. Ces deux approches de diagnostic ont été construites pour pouvoir être applicable en ligne et ainsi garantir un suivi non intrusif du vieillissement des systèmes de stockage.

Au-delà du diagnostic instantané, la problématique de l'anticipation a été traitée par le développement d'outils de pronostic. L'utilisation de réseaux de neurones LSTM a notamment permis une approche généralisable pour prédire l'évolution des états de santé ainsi que du temps de vie restant des composants. Cette approche a été complétée par une quantification des incertitudes de prédiction basée sur le Monte-Carlo \textit{dropout}.

Enfin, l'intégration de ces différents indicateurs au sein d'un jumeau numérique a été présentée comme un moyen d'assurer une cohérence entre le système réel et sa représentation virtuelle, facilitant ainsi la détection de dérives anormales.

L'obtention de ces informations de santé ne constitue toutefois qu'une étape préliminaire. La valeur de ces travaux réside dans l'exploitation de ces données pour orienter les décisions de pilotage. Pour garantir notamment la viabilité et la durabilité du micro-réseau, la connaissance de la dégradation doit être intégrée à la couche de commande du système.
Ce changement de paradigme guide les développements du chapitre suivant, consacré à la construction d'un cadre méthodologique pour la stratégie de gestion d'énergie. Les critères et les indicateurs de performance nécessaires à l'évaluation objective des commandes y sont d'abord définis. Par la suite, plusieurs stratégies de gestion d'énergie de référence sont présentées. Celles-ci servent de socle comparatif et de base méthodologique aux développements originaux proposés dans la suite du manuscrit. Les connaissances du vieillissement sont intégrées dans la stratégie de pilotage, avec pour objectif une gestion de l'énergie qui allie sécurité de l'approvisionnement et soutenabilité dans le temps.